% DRAFT — KJV context, for Ch 1 or Ch 2 near the first scripture evidence
% Explains WHY the Bible is the example, connecting to Bloom

The choice of the King James Bible as the first experimental subject is not arbitrary. Harold Bloom argued throughout his career that the KJV stands alongside Shakespeare as a constitutive text of the English language---not a translation of Hebrew and Greek originals but a new literary work so powerful that it became a source of the language's imaginative possibilities.\footnote{Harold Bloom, \emph{The Shadow of a Great Rock: A Literary Appreciation of the King James Bible}, Yale University Press, 2011. Bloom's specific claim is that the KJV is a ``strong reading'' of its sources so thorough that it ``constitutes a new original.''} If the manifold is a compressed portrait of a civilisation's writing, then the KJV is among the deepest strata of the English-language manifold---a text whose phrases, cadences, and imagery have been training English speakers for four centuries before any language model existed. To embed the KJV and trace its trajectory through meaning-space is to watch one of the geological layers of the manifold reveal its internal structure: the basins, returns, and ruptures that were deposited into the geometry of English by a committee of translators whose stylistic choices still shape what counts as authoritative, sacred, and beautiful in the language.

The KJV is also an evolving text in the book's technical sense. Read sequentially---Genesis through Revelation---each verse carries the accumulated context of everything that precedes it. The trajectory is not a static map but a dynamic traversal: what ``covenant'' means in Exodus is different from what it means in Romans, because the attention field has changed, because the trajectory has visited territory that reshapes the word's coordinates. The embedding model reads the Bible the way a careful reader does: carrying forward what has been said, letting it inflect what comes next.
